\documentclass{article}
\usepackage[T1]{fontenc}
\usepackage{lmodern}
\usepackage{graphicx}
\usepackage{amsmath,amssymb}
\usepackage[a4paper,margin=2cm]{geometry}
\usepackage[absolute,overlay]{textpos}
\setlength{\TPHorizModule}{1mm}
\setlength{\TPVertModule}{1mm}
\usepackage[indonesian]{babel}
\usepackage{hyperref}
\setlength{\emergencystretch}{2em}
% unit: Halaman 2

\begin{document}

% === Halaman 2 (python/native) ===

2

Pendahuluan

• Kriptografi klasik (classical cryptography) merupakan kriptografi yang sudah tua, 

sudah ada sejak ribuan tahun yang lalu hingga ditemukan komputer digital

• Cipher pada kriptografi klasik, dinamakan cipher klasik, (classical cipher) hanya 

memproses pesan dari huruf-huruf alfabet saja 

• Menggunakan alat tulis pena dan kertas saja, belum ada komputer

• Termasuk ke dalam jenis kriptografi kunci-simetri

• Tiga alasan mempelajari kriptografi klasik:

 1. Memahami konsep dasar kriptografi.

 2. Sebagai dasar algoritma kriptografi modern. 

 3. Untuk memahami kelemahan sistem cipher.

\end{document}
